% s02-curve-a.tex —— §1.1–§1.4（原文 p.3–p.10）
% 底本：L. Fargues, Geometrization of the local Langlands correspondence: an overview, 2016-02-02
% 体例：每页起点 \srcpage{N}；定理/公式编号与原文逐条一致；译注 \yizhu{注文}{来源}。
% 注：本文件不写 documentclass / \begin{document}。
% ============================================================

% 原文用 ⟦⟧ 表示 π-进圆括号（如 R°⟦π⟧）；本工程 preamble 未载 stmaryrd，
% 故就地给出兜底定义（若他处已定义则不覆盖）。

\section{曲线}% 原文 1. The curve
\setcounter{equation}{0}% 原文公式 (1)–(5) 自本页起编号，译文保号

\subsection{adic 曲线}% 原文 1.1. The adic curve

\leavevmode\srcpage{3}设 $E$ 是剩余域有限的局部域，剩余域为 $\FF_q=\oo_E/\pi$。于是或者 $[E:\QQ_p]<+\infty$，或者 $E=\FF_q((\pi))$。我们记
\[
  \mathrm{Perf}_{\FF_q}
\]
为 $\FF_q$ 上完满空间构成的范畴\yizhu{「完满空间」不固定底域：Scholze 的原始定义要求底域是完满域，本文采用 [11]、[27] 中允许任意基的推广；否则「$S$ 上的一族曲线」无从谈起。}{A：原文 §1.1 p.3（引 [11]、[27]）；I：本工程 glossary.tsv}。这里只须说明：在本文中，完满空间并不像 Scholze 原始定义（[24]）那样生活在一个固定的基完满域之上。我们采用 [11] 或 [27] 中更一般的定义。

给定 $S\in\mathrm{Perf}_{\FF_q}$，便可定义一个 $E$-adic 空间 $X_S$\yizhu{「adic 空间」是 Huber 引入的框架：它同时带经典点与几何点，比刚性解析空间多出「边界」上的点；本文的曲线 $X_S$ 就住在这个范畴里。}{A：Huber, \emph{Étale cohomology of rigid analytic varieties}（1996）；I：本工程 glossary.tsv}：按定义，它被一个空间 $Y_S$ 单值化，
\[
  X_S=Y_S/\varphi^\ZZ,
\]
其中 $\varphi$ 是一个在 $Y_S$ 上真不连续作用的 Frobenius 作用。

\subsubsection*{1.1.1. 等特征情形}% 原文 1.1.1. The equal characteristic case

$E=\FF_q((\pi))$ 的情形更简单。事实上我们有下面的定义。

\srcpage{4}\begin{definition}\label{def:1.1}
若 $E=\FF_q((\pi))$，令 $Y_S=\mathbb{D}^*_S\subset\AA^1_S$ 为变量 $\pi$ 的一个去心开圆盘，其中 $\mathbb{D}^*_S=\{0<|\pi|<1\}$。作用在 $\mathbb{D}^*_S$ 上的 Frobenius $\varphi$ 就是 $S$ 的那个 Frobenius。
\end{definition}

若 $S=\Spa(R,R^+)$ 是 affinoid 完满的，则
\[
  \mathbb{D}^*_S=\Spa\bigl(R^\circ\llbracket\pi\rrbracket,\;R^++\pi R^\circ\llbracket\pi\rrbracket\bigr)\setminus V(\pi\varpi_R),
\]
其中 $\varpi_R\in R^{\circ\circ}\cap R^\times$ 就是所谓的伪一致化元（[11]）。此外，$\varphi$ 由
\[
  \varphi\Bigl(\sum_{n\ge0}x_n\pi^n\Bigr)=\sum_{n\ge0}x_n^q\pi^n
\]
作为 $R^\circ\llbracket\pi\rrbracket$ 的自同构诱导。$\varpi_R$ 的选取给出一个半径函数
\begin{equation}\label{eq:1.1}
\begin{aligned}
  \delta:\ \bigl|\mathbb{D}^*_S\bigr| &\longrightarrow\ ]0,1[,\\
  y &\longmapsto p^{-\frac{\log|\pi(y^{\max})|}{\log|\varpi_R(y^{\max})|}}\;,
\end{aligned}
\end{equation}
\yizhu{$\delta(y)$ 度量 $y$ 处 $\pi$ 相对 $\varpi_R$ 的 $p$-进「半径」；$\delta(\varphi(y))=\delta(y)^{1/q}$ 说明 $\varphi$ 把半径开 $q$ 次方，故 $\varphi^\ZZ$ 自由、真不连续。}{A：原文 §1.1.1 p.4}其中 $y^{\max}$ 是 $y$ 的极大泛化，即一个其值群可取为 $\RR$ 之子群的赋值（$y^{\max}$ 是通过把 Berkovich 谱等同于 adic 谱的商而联系于 $y$ 的 Berkovich 点）。这个连续函数满足
\[
  \delta(\varphi(y))=\delta(y)^{1/q},
\]
于是 $\varphi$ 的作用真不连续且没有不动点。此时有两个结构态射
\begin{center}
\begin{tikzcd}
  & \mathbb{D}^*_S \arrow[dl] \arrow[dr] & \\
  S & & \mathbb{D}^*_{\FF_q}=\Spa(E)
\end{tikzcd}
\end{center}
其中：

\begin{itemize}
  \item 第一个是去心圆盘的标准结构态射：它局部有限型，但\emph{不是} $\varphi$-不变的。
  \item \srcpage{5}第二个通过把 $\FF_q$ 上的去心开圆盘（赋以平凡赋值）与 $\Spa(E)$ 等同而得到。它 $\varphi$-不变，但\emph{不是}局部有限型。这才是我们关心的结构态射。
\end{itemize}

\begin{definition}\label{def:1.2}
令
\[
  X_S=\mathbb{D}^*_S/\varphi^\ZZ
\]
为一个 $E$-adic 空间。
\end{definition}

这个空间 $X_S$ 不是 $E$ 上的有限型。而且它没有到 $S$ 的结构态射。尽管如此，对 $\tau\in\{\text{analytic},\text{étale},\text{pro-étale}\}$，函子 $T/S\mapsto X_T/X_S$ 定义了站点间的连续态射
\[
  (X_S)_\tau\longrightarrow S_\tau.
\]
可以这样理解这件事：即使空间「$S/\mathrm{Frob}^\ZZ$」没有意义，与之相伴的站点「$(S/\mathrm{Frob}^\ZZ)_\tau$」也必定与 $S_\tau$ 重合。特别地，有连续映射
\[
  |X_S|\longrightarrow|S|,
\]
它拟紧且特殊化（$f$ 特殊化，指只要 $f(x)\succ y$，就存在 $x\succ x'$ 使 $f(x')=y$）；因此当 $S$ 拟紧拟分离时，它还是闭映射（谱空间之间的谱映射是闭映射，当且仅当它是特殊化映射）。于是可以把 $X_S$ 想成「在 $S$ 上紧合」。

当 $S$ 是完满域 $F$ 的谱时，空间 $X_F$ 具有很强的有限性：
\begin{itemize}
  \item 它在 Huber 的意义下强 Noether；
  \item 它「是一条曲线」：它被有限个 $E$-Banach 代数的谱覆盖，而这些代数都是主理想整环。
\end{itemize}
这些性质是清楚的：把 $\mathbb{D}^*_F$ 看成 $F$ 上的经典 Tate 刚性空间时，闭圆环上全纯函数构成的代数都是主理想整环。但对一般的 $S$ 就不再如此：一般 $S$ 上的空间 $X_S$ 不满足任何合理的有限性。话虽如此，仍可把 $X_S$ 想成曲线族
\[
  (X_{k(s)})_{s\in S},
\]
尽管无法给这句话一个精确的含义。

\subsubsection*{1.1.2. 不等特征情形}% 原文 1.1.2. The unequal characteristic case

现在设 $[E:\QQ_p]<+\infty$。令 $S=\Spa(R,R^+)$ 为 affinoid 完满空间。前面那个环 $R^\circ\llbracket\pi\rrbracket$ 的类比物是 Fontaine 的环 $A_{\mathrm{inf}}$，记作 $\mathbf{A}$。

\begin{definition}\label{def:1.3}
定义
\[
\begin{aligned}
  \mathbf{A} &= W_{\oo_E}(R^\circ)\\
  &= W(R^\circ)\otimes_{W(\FF_q)}\oo_E\\
  &= \Bigl\{\sum_{n\ge0}[x_n]\pi^n \;\Big|\; x_n\in R^\circ\Bigr\}\qquad\text{（\emph{唯一}写法）},
\end{aligned}
\]
其中记号 $W_{\oo_E}$ 指分歧 Witt 向量。令
\[
  \mathbf{A}^+=\Bigl\{\sum_{n\ge0}[x_n]\pi^n\in\mathbf{A}\;\Big|\;x_0\in R^+\Bigr\}.
\]
\end{definition}

\begin{definition}\label{def:1.4}
给 $\mathbf{A}$ 赋以 $(\pi,[\varpi_R])$-进拓扑，并令
\[
  Y_{R,R^+}=\Spa(\mathbf{A},\mathbf{A}^+)\setminus V(\pi[\varpi_R]).
\]
\end{definition}

在公式 \eqref{eq:1.1} 中把 $\varpi_R$ 换成 $[\varpi_R]$，就得到连续函数
\[
  \delta:\ |Y_S|\longrightarrow\ ]0,1[.
\]
固定一个幂乘性 Banach 代数范数 $\|\cdot\|:R\to\RR_+$，它定义 $R$ 的拓扑（一旦 $\|\varpi_R\|$ 固定，这样的范数唯一）。空间 $Y_{R,R^+}$ 是 \emph{Stein} 的，因而完全由 $\cO(Y_{R,R^+})$ 决定。后者可这样描述。考虑
\[
\begin{aligned}
  \cO(Y_{R,R^+})^b &= \mathbf{A}\Bigl[\tfrac{1}{\varpi_E},\tfrac{1}{[\varpi_R]}\Bigr]\\
  &= \Bigl\{\sum_{n\gg-\infty}[x_n]\pi^n \;\Big|\; x_n\in R,\ \sup_n\|x_n\|<+\infty\Bigr\},
\end{aligned}
\]
即 $Y_{R,R^+}$ 上沿除子 $\pi=0$ 与 $[\varpi_F]=0$ 亚纯的全纯函数构成的环\yizhuw{原文此处作 $[\varpi_F]=0$；按上下文（商环 $\mathbf{A}[\frac{1}{\varpi_E},\frac{1}{[\varpi_R]}]$）应为 $[\varpi_R]=0$。译文照原样。}。对每个 $\rho\in]0,1[$，$\cO(Y_{R,R^+})^b$ 上有一个幂乘性 Gauss 范数 $\|\cdot\|_\rho$，定义为
\[
  \Bigl\|\sum_{n\gg-\infty}[x_n]\pi^n\Bigr\|_\rho=\sup_{n\in\ZZ}\|x_n\|\rho^n.
\]
于是 $\cO(Y_{R,R^+})$ 就是 $\cO(Y_{R,R^+})^b$ 关于 $(\|\cdot\|_\rho)_{\rho\in]0,1[}$ 完备化所得的 $E$-Fréchet 代数。

Frobenius $\varphi$ 在 $Y_{R,R^+}$ 上的作用由 Witt 向量通常的 Frobenius 给出
\[
  \varphi\Bigl(\sum_{n\ge0}[x_n]\pi^n\Bigr)=\sum_{n\ge0}[x_n^q]\pi^n.
\]
然后可以验证，函子 $(R,R^+)\mapsto Y_{R,R^+}$ 粘成函子 $S\mapsto Y_S$，并定义 $X_S=Y_S/\varphi^\ZZ$。此外 §1.1.2 的所有性质对 $X_S$ 都成立\yizhuw{原文此处作「§1.1.2 的所有性质」；本节即 §1.1.2，所指应是上一小节 §1.1.1 中列出的性质。译文照原样。}（尽管其中一些性质难证得多，典型的是：对完满域 $F$，$X_F$「是一条强 Noether 曲线」，见 [10] 与 [15]）。

\subsection{等特征与不等特征的联系}% 原文 1.2. The link between the equal and unequal characteristic case

\leavevmode\srcpage{6}事实上，前面的 $E$-adic 空间 $Y_{S,E}$ 与 $X_{S,E}=Y_{S,E}/\varphi^\ZZ$ 对任何剩余域为 $\FF_q$ 的完备非阿基米德域 $E$ 都可定义。我们只关心 $E$ 有离散赋值（前述情形）或者 $E$ 完满的情形\yizhuw{原文此处作「(the preceding case) of $E$ is perfectoid」，其中 of 当为 or 之误；译文按 or 读作「或者 $E$ 是完满的」。}。事实是，这个构造与基变换兼容：
\[
  Y_{S,E'}=Y_{S,E}\hat\otimes_EE',\qquad X_{S,E'}=X_{S,E}\hat\otimes_EE'.
\]
当然，当 $E|\QQ_p$ 时，这些公式把 $Y_{S,E}$ 与 $X_{S,E}$ 完全化归为 $E$ 有离散赋值的情形：取 $E_0=W(\FF_q)\bigl[\tfrac1p\bigr]$，它们给出
\[
  Y_{S,E}=Y_{S,E_0}\hat\otimes_{E_0}E,\qquad X_{S,E}=X_{S,E_0}\hat\otimes_{E_0}E.
\]
当 $E$ 有不等特征时就不再如此，因为先验地，那些公式给出的空间定义可能依赖于 $E$ 内子离散赋值域的选取。我们要说的是：对任何 $E$，仍有自然的定义。确切地说，下面是定义兼命题。

\begin{definition}\label{def:1.5}
设 $E$ 的剩余域为 $\FF_q$，但不必有离散赋值。
\begin{enumerate}[label=(\arabic*)]
  \item 若 $E|\QQ_p$，令 $Y_{S,E}=Y_{S,E_0}\hat\otimes_{E_0}E$，$X_{S,E}$ 同此。
  \item 对剩余域为 $\FF_q$、特征为 $p$ 的 $E$，以及 $\FF_q$ 上的 affinoid 完满空间 $(R,R^+)$，令
  \[
    Y_{R,R^+}=\Spa(R^\circ\hat\otimes_{\FF_q}\oo_E)\setminus V(\varpi_E\varpi_R),
  \]
  其中 $\varpi_E$ 是 $E$ 中的伪一致化元。这粘成函子 $S\mapsto Y_{S,E}$。
\end{enumerate}
\end{definition}

事实上，当 $E$ 有特征 $p$ 时有一个简单的公式（其证明相当初等）：
\begin{equation}\label{eq:1.2}
  Y_{S,E}=S\times_{\Spa(\FF_q)}\Spa(E),
\end{equation}
作为 $\FF_q$ 上（解析）adic 空间范畴中的纤维积。这个公式虽然简单而优雅，却没有说出 $Y_{S,E}$ 的任何几何（典型的问题：它上面的向量丛是什么？）。注意通过这个公式
\[
  \varphi=\mathrm{Frob}_S\times\mathrm{Id},
\]
于是我们得到简单的公式
\begin{equation}\label{eq:1.3}
  X_{S,E}=(S\times_{\Spa(\FF_q)}\Spa(E))/\mathrm{Frob}_S^\ZZ\times\mathrm{Id},
\end{equation}
其中的商是解析拓扑下的层商。

例如设 $E=\FF_q((\pi^{1/p^\infty}))$。则
\[
  Y_{S,E}=\varprojlim_{\mathrm{Frob}}\mathbb{D}^*_S=\mathbb{D}^*_S\hat\otimes_{\FF_q((\pi))}\FF_q((\pi^{1/p^\infty})),
\]
即完满的去心开圆盘。

我们刚刚指出，当 $E=\FF_q((\pi^{1/p^\infty}))$ 时，空间 $Y_{S,E}$、从而 $X_{S,E}$ 是完满的。这在等特征与不等特征中其实都更一般。

\begin{proposition}\label{prop:1.6}
对任何完满的 $E$，空间 $Y_{S,E}$ 与 $X_{S,E}$ 都完满。对这种完满的 $E$，还有下述非常简单的公式
\[
  Y^\flat_{S,E}=Y_{S,E^\flat},\qquad X^\flat_{S,E}=X_{S,E^\flat}.
\]
\end{proposition}

\begin{remark}\label{rem:1.7}
对离散赋值的 $E$，$Y_{(R,R^+),E}$ 在标量扩张到一个完满域之后变成完满的；这一事实是证明 $Y_{(R,R^+),E}$ 为 adic 空间（即 Huber 的预层是层）的要点。这本应在进入本节之前就指出，但既然这是一篇综述，作者选择少管严格的严谨性，多顾阐述。
\end{remark}

\leavevmode\srcpage{7}让我们处理一个例子。取 $[E:\QQ_p]<+\infty$。设 $E_\infty|E$ 是由 $\oo_E$ 上一条 Lubin–Tate 形式群律 $\mathcal{L}T$ 的挠点所生成扩域的完备化。Lubin–Tate 特征诱导同构 $\Gal(E_\infty|E)\xrightarrow{\sim}\oo_E^\times$。于是有
\[
  E_\infty^\flat=\FF_q((T^{1/p^\infty})),
\]
其中若 $\varepsilon=(\varepsilon^{(n)})_{n\ge0}$ 是 $\mathcal{L}T$ 的 Tate 模的一组生成元，$[\pi]_{\mathcal{L}T}(\varepsilon^{(n+1)})=\varepsilon^{(n)}$，则
\[
  T=(\varepsilon^{(n)}\bmod\pi)_{n\ge0}\in\varprojlim_{\mathrm{Frob}}\oo_{E_\infty}/\pi=\oo^\flat_{E_\infty}.
\]
元素 $a\in\oo_E^\times=\Gal(E_\infty|E)$ 在 $E_\infty^{\flat,\times}$ 上的作用由 $[a]_{\mathcal{L}T}\in\Aut(\FF_q\llbracket T\rrbracket)$ 给出。作用 $[-]_{\mathcal{L}T}$ 诱导 $\oo_E^\times$ 在 $\mathbb{D}^*_S$ 上的作用，并且有 $\oo_E^\times$-等变的等同
\[
  \bigl(Y_{S,E}\hat\otimes_EE_\infty\bigr)^\flat=Y^\flat_{S,E_\infty}=Y_{S,E_\infty^\flat}=\varprojlim_{\mathrm{Frob}}\mathbb{D}^*_S.
\]

\subsection{菱形乘积公式}% 原文 1.3. The diamond product formula

我们将说明，公式 \eqref{eq:1.2} 借助 Scholze 的菱形理论（[27]）推广到不等特征。下文还会系统地使用菱形，因此这里先对该理论作一快速回顾。

\subsubsection*{1.3.1. 菱形}% 原文 1.3.1. Diamonds

考虑 $[E:\QQ_p]<+\infty$ 的情形。我们刚看到，虽然 $Y_{F,E}$ 不完满，$Y_{F,E}\hat\otimes E_\infty$ 却完满，其倾斜为 $\mathbb{D}^{*,1/p^\infty}_S$，即完满去心圆盘。于是人们忍不住「下降」，说 $Y_{S,E}$ 就是
\begin{equation}\label{eq:1.4}
  \text{“}\mathbb{D}^{*,1/p^\infty}_S/\oo_E^\times\text{”}.
\end{equation}
这概括在下面的\emph{菱形}交换图中\yizhu{「菱形」是 pro-étale 拓扑下的代数空间，即完满空间按 pro-étale 等价关系取的商。有了它，像 $Y_{S,E}=\mathbb{D}^{*,1/p^\infty}_S/\oo_E^\times$ 这样的写法才有意义。}{A：原文 §1.3.1 p.7 与 Definition 1.12；I：本工程 glossary.tsv}：
\begin{center}
\begin{tikzcd}
  & Y_{S,E_\infty} \arrow[dl, "(-)^\flat"'] & \\
  \mathbb{D}^{*,1/p^\infty}_S \arrow[dr] & & Y_{S,E} \arrow[ul, "-\hat{\otimes}E_\infty"] \arrow[dl, "="'] \\
  & \mathbb{D}^{*,1/p^\infty}_S/\oo_E^\times &
\end{tikzcd}
\end{center}

这是一个更一般的现象。典型地，考虑环面 $\Spa(\QQ_p\langle T^{\pm1}\rangle)$。它有一个 pro-Galois 的完满覆盖
\[
  \Spa(\QQ_p^{\mathrm{cyc}}\langle T^{\pm1/p^\infty}\rangle)\longrightarrow\Spa(\QQ_p\langle T^{\pm1}\rangle).
\]
与前面同样的论证使我们写出
\begin{equation}\label{eq:1.5}
  \text{“}\Spa(\QQ_p\langle T^{\pm1}\rangle)=\Spa(\QQ_p^{\mathrm{cyc},\flat}\langle T^{\pm1/p^\infty}\rangle)/\ZZ_p(1)\rtimes\Gal(\QQ_p^{\mathrm{cyc}}|\QQ_p)\text{”}.
\end{equation}
菱形理论让我们能给公式 \eqref{eq:1.4} 与 \eqref{eq:1.5} 赋予意义。

为此，给 $\mathrm{Perf}_{\FF_q}$ 与 $\mathrm{Perf}_E$ 赋以 pro-étale 拓扑。

\begin{definition}\label{def:1.8}
\begin{enumerate}[label=(\arabic*)]
  \item 若 $S\in\mathrm{Perf}_{\FF_q}$，$S$ 的一个 \emph{untilt} 按定义是一对 $(S^\sharp,\iota)$，其中 $S^\sharp\in\mathrm{Perf}_E$，$\iota\colon S\xrightarrow{\sim}S^{\sharp,\flat}$\yizhu{untilt 与倾斜互逆：倾斜把特征 $0$ 的完满空间送到特征 $p$；untilt 则是给定特征 $p$ 的 $S$，配上它在特征 $0$ 中的提升 $S^\sharp$ 与同构 $\iota$。此词不译。}{A：原文 Definition 1.8 p.7；I：本工程 glossary.tsv}。
  \item 若 $\mathcal{F}$ 是 $\mathrm{Perf}_E$ 上的 pro-étale 层，定义 $\mathcal{F}^\diamond$ 为 $\mathrm{Perf}_{\FF_q}$ 上的层
  \[
    \mathcal{F}^\diamond(S)=\bigl\{(S^\sharp,\iota,s)\ \big|\ (S^\sharp,\iota)\ \text{是}\ S\ \text{的 untilt，且}\ s\in\mathcal{F}(S^\sharp)\bigr\}/\sim.
  \]
\end{enumerate}
\end{definition}

\begin{remark}\label{rem:1.9}
$\mathcal{F}^\diamond$ 是层，这一点依赖 Scholze 的纯度。事实上，它可由下述事实推出：倾斜诱导 pro-étale 站点的等价 $\mathrm{Perf}_{S^\sharp}\xrightarrow{\sim}\mathrm{Perf}_S$。
\end{remark}

\leavevmode\srcpage{8}例如 $\Spa(E)^\diamond$ 就是 $E$ 的 untilt 构成的层（模去显然的等价关系）。多数时候为了简化记号，我们会说「$S^\sharp$ 是 $S$ 的一个 untilt」，尽管应该说：我们考虑一对 $(S^\sharp,\iota)$，模去显然的等价关系。

前述菱形函子给出站点与意象的等价
\[
  \mathrm{Perf}_E\xrightarrow{\sim}\mathrm{Perf}_{\FF_q}/\Spa(E)^\diamond
\]
以及
\[
  \widetilde{\mathrm{Perf}_E}\xrightarrow{\sim}\widetilde{\mathrm{Perf}_{\FF_q}/\Spa(E)^\diamond}.
\]
当然，若 $\mathcal{F}$ 可由某个 $T\in\mathrm{Perf}_E$ 表示，则 $T^\diamond=T^\flat$。

于是我们遇到下面的问题：若 $X$ 是 $E$ 上的刚性空间，它自然地定义 $\mathrm{Perf}_E$ 上的一个层。看这个层会不会丢掉 $X$ 的某些性质？换个说法：$X^\diamond/\Spa(E)^\diamond$ 是否刻画 $X$？

对正规的刚性空间，答案是肯定的，理由是下面两条性质。

\begin{proposition}\label{prop:1.10}
\begin{enumerate}[label=(\arabic*)]
  \item 任何刚性解析空间都有完满的 pro-étale 覆盖。
  \item 若 $X$ 还是正规的，则对相伴的站点态射 $\nu:X_{\mathrm{pro\text{-}ét}}\to X_{\mathrm{ét}}$ 有 $\nu_*\wt{\cO}_X=\cO_X$，其中 $\wt{\cO}_X(R,R^+)=R$（当 $\Spa(R,R^+)\to X$ 是 pro-étale 的 affinoid 完满映射时）（[26]）。
\end{enumerate}
\end{proposition}

为说明这些性质，考虑 $X=\Spa(E\langle T_1^{\pm1},\dots,T_d^{\pm1}\rangle)$。则
\[
  \Spa(C_p\langle T_1^{\pm1/p^\infty},\dots,T_d^{\pm1/p^\infty}\rangle)\longrightarrow X
\]
是一个完满的 pro-étale Galois 覆盖，Galois 群为 $\ZZ_p(1)^d\rtimes\Gal(\ol{E}|E)$。这就对 $E$ 上的光滑刚性空间证明了性质 (1)。再经显式计算可验证
\[
  \bigl[C_p\langle T_1^{\pm1/p^\infty},\dots,T_d^{\pm1/p^\infty}\rangle\bigr]^{\ZZ_p(1)^d\rtimes\Gal(\ol{E}|E)}=E\langle T_1^{\pm1},\dots,T_d^{\pm1}\rangle.
\]
这就对 $E$ 上的光滑刚性空间证明了性质 (2)。

\begin{corollary}\label{cor:1.11}
我们有 $E$ 上正规刚性空间范畴到 $\widetilde{\mathrm{Perf}}_{\FF_q}/\Spa(E)^\diamond$ 的嵌入。
\end{corollary}

但事实上，函子 $X\mapsto X^\diamond$ 的本质像具有特殊的类型，例如从公式 \eqref{eq:1.5} 可以看出：它是 $\mathrm{Perf}_{\FF_q}$ 中 pro-étale 拓扑下的代数空间。这正是菱形的定义。

\begin{definition}[Scholze {[27]}]\label{def:1.12}
\emph{菱形}是 $\mathrm{Perf}_{\FF_p}$ 中 pro-étale 拓扑下的代数空间。
\end{definition}

若 $X$ 是 $E$ 上的正规刚性空间，取一个完满的 pro-étale 覆盖 $\wt{X}\to X$。可以验证（这需要一点工作，但并不很难）：
\begin{itemize}
  \item 纤维积 $\mathcal{R}=\wt{X}\times_X\wt{X}$ 与 $\wt{X}\times_{\Spa(E)}\wt{X}$ 作为完满空间存在，且 $(\wt{X}\times_{\Spa(E)}\wt{X})^\flat=\wt{X}^\flat\times_{\Spa(\FF_q)}\wt{X}^\flat$。
  \item 投影 $\mathcal{R}\rightrightarrows\wt{X}$ 是 pro-étale 的。
\end{itemize}
于是 $\mathcal{R}^\flat$ 是 $\wt{X}^\flat$ 上的 pro-étale 等价关系，并且
\[
  X^\diamond=\wt{X}^\flat/\mathcal{R}^\flat.
\]
于是可以陈述下面的结果。

\begin{corollary}\label{cor:1.13}
有完全忠实函子
\[
\begin{array}{rcl}
  E\ \text{上的正规刚性空间} &\longrightarrow& E^\diamond\ \text{上的菱形}\\[2pt]
  X &\longmapsto& X^\diamond/E^\diamond.
\end{array}
\]
\end{corollary}

\subsubsection*{1.3.2. 乘积公式}% 原文 1.3.2. The product formula

\leavevmode\srcpage{9}同样的原则适用于 $E$ 上的前完满空间：即在任意完满扩域上作扩张后成为完满的 $E$-adic 空间。例如 $\Spa(E\langle T^{\pm1/p^\infty}\rangle)$ 或 $Y_{S,E}$ 都是。现在可以计算 $Y^\diamond_{S,E}$\yizhu{「前完满」指在任意完满扩域上作基变换后成为完满的 $E$-adic 空间，自身未必完满；$Y_{S,E}$ 就是这种情形，所以要先取菱形 $Y^\diamond_{S,E}$ 才能谈「$S$ 上的空间」。}{A：原文 §1.3.2 p.9}：
\[
\begin{aligned}
  Y^\diamond_{S,E} &= Y_{S,E_\infty^\flat}/\oo_E^\times\\
  &= S\times\Spa(E_\infty^\flat)/\oo_E^\times\\
  &= S\times\Spa(E)^\diamond.
\end{aligned}
\]
我们于是找到与 \eqref{eq:1.2} 类似的公式。最终有
\[
  X^\diamond_{S,E}=(S\times\Spa(E)^\diamond)/\varphi^\ZZ,
\]
其中 $\varphi=\mathrm{Frob}_S\times\mathrm{Id}$。

\begin{remark}\label{rem:1.14}
$Y^\diamond_{S,E}$ 的计算其实可以不用辅助扩域 $E_\infty$。稍微深挖一下就会发现，这是下述事实的推论：有两个伴随函子
\[
\begin{tikzcd}
  \text{完满}\ \FF_p\text{-代数} \arrow[r, shift left=0.7ex, "W(-)"] & p\text{-进环} \arrow[l, shift left=0.7ex, "(-)^\flat"]
\end{tikzcd}
\]
其中伴随态射是 Fontaine 的 theta $\theta:W(A^\flat)\to A$ 与 $R\xrightarrow{\sim}W(R)^\flat$，后者由 $x\mapsto\bigl([x^{p^{-n}}]\bigr)_{n\ge0}$ 给出。
\end{remark}

\subsection{untilt 作为曲线上的 Cartier 除子}% 原文 1.4. Untilts as Cartier divisors on the curve

在 [10] 中，作者们证明了：他们构造的曲线上，与完满域 $F|\FF_q$ 相伴的闭点对应 $F$ 的 untilt 模 Frobenius 的幂。这一点在相对情形下推广。

\subsubsection*{1.4.1. 等特征情形}% 原文 1.4.1. The equal characteristic case

设 $E=\FF_q((\pi))$。取 $S\in\mathrm{Perf}_{\FF_q}$。由于 $Y_S=S\times\Spa(E)$，任何态射 $S\to\Spa(E)$ 都诱导结构态射 $Y_S\to\Spa(E)$ 的一个截面。事实上这是一个闭嵌入
\[
  S\hookrightarrow Y_S=\mathbb{D}^*_S,
\]
当 $S=\Spa(R,R^+)$ 时，它局部由 $(\pi-a)$ 定义，其中 $a\in R^{\circ\circ}\cap R^\times$。人们把它想成 $Y_S$ 上一个次数 $1$ 的 Cartier 除子。

\begin{definition}\label{def:1.15}
元素 $f=\sum_{n\ge0}a_n\pi^n\in R^\circ\llbracket\pi\rrbracket$ 称为\emph{次数 $1$ 本原的}，若 $a_0\in R^{\circ\circ}\cap R^\times$ 且 $a_1\in R^{\circ\times}$。
\end{definition}

根据 Weierstrass 分解\yizhu{「次数 $1$ 本原元」是 Weierstrass 意义下首项可逆的一次多项式：等特征时它总能写成单位乘 $\pi-a$，因而对应 $S$ 到 $\Spa(E)$ 的一个截面，也就是 $Y_S$ 上次数 $1$ 的 Cartier 除子。}{A：原文 §1.4.1 p.9}，$R^\circ\llbracket\pi\rrbracket$ 中任何次数 $1$ 的本原元都可以写成一个单位乘以 $\pi-a$，其中 $a\in R^{\circ\circ}\cap R^\times$ 由该元唯一确定。由此得到下面的引理。

\begin{lemma}\label{lem:1.16}
有双射
\[
\begin{aligned}
  &\{\text{态射}\ S\to\Spa(E)\}\\
  &\simeq\ \{\text{闭浸入}\ T\hookrightarrow Y_S,\ \text{局部由次数 }1\text{ 本原元定义}\}.
\end{aligned}
\]
\end{lemma}

\subsubsection*{1.4.2. 不等特征情形}% 原文 1.4.2. The unequal characteristic case

现在设 $[E:\QQ_p]<+\infty$。令 $R$ 是完满的 $\FF_q$-代数。

\begin{definition}\label{def:1.17}
$f=\sum_{n\ge0}[a_n]\pi^n\in\mathbf{A}=W_{\oo_E}(R^\circ)$ 称为\emph{次数 $1$ 本原元}，若 $a_0\in R^{\circ\circ}\cap R^\times$ 且 $a_1\in R^{\circ\times}$。
\end{definition}

前面关于 Weierstrass 分解的陈述这时不再成立。不过我们有下面的结果。

\begin{proposition}\label{prop:1.18}
\begin{enumerate}[label=(\arabic*)]
  \item 若 $R^\sharp$ 是 $R$ 在 $E$ 上的一个 untilt，则 Fontaine 的 $\theta:W_{\oo_E}(R^\circ)\twoheadrightarrow R^{\sharp,\circ}$ 的核由一个次数 $1$ 本原元生成。
  \item 反之，若 $f\in W_{\oo_E}(R^\circ)$ 是本原的，则 $W_{\oo_E}(R^\circ)\bigl[\tfrac1\pi\bigr]/f$ 是一个完满的 $E$-代数，且是 $R$ 的一个 untilt。
\end{enumerate}
\end{proposition}

\srcpage{10}\begin{example}\label{ex:1.19}
取 $F=\QQ_p^{\mathrm{cyc},\flat}=\FF_p((T^{1/p^\infty}))$，其中 $T=\bigl((\zeta_{p^n}-1)\bmod p\bigr)_{n\ge0}\in\varprojlim_{\mathrm{Frob}}\ZZ_p^{\mathrm{cyc}}/p=\ZZ_p^{\mathrm{cyc},\flat}$。则 $\ZZ_p^{\mathrm{cyc}}=W(\FF_p\llbracket T^{1/p^\infty}\rrbracket)\big/\bigl(1+[1+T^{1/p}]+\cdots+[1+T^{\frac{p-1}{p}}]\bigr)$。此时本原元 $1+[1+T^{1/p}]+\cdots+[1+T^{\frac{p-1}{p}}]$ 没有 Weierstrass 分解，即不能写成一个单位乘以 $p-[a]$（对某个 $a$），因为 $p$ 在 $\QQ_p^{\mathrm{cyc}}$ 中没有 $p$ 次方根。
\end{example}

\begin{remark}\label{rem:1.20}
若 $R^\sharp=W_{\oo_E}(R^\circ)\bigl[\tfrac1\pi\bigr]/f$，则在 $R$ 的一个 pro-étale 覆盖之后，本原元 $f$ 有 Weierstrass 分解，即写成一个单位乘以 $\pi-[a]$。事实上，只需对 $R^\sharp$ 由 $\bigcup_{n\ge0}R^\sharp(\pi^{1/p^n})$ 的完备化给出的 pro-étale 覆盖取倾斜。不过这样的 $a$ 不像等特征情形那样唯一。
\end{remark}

由此得到下面的推论。

\begin{corollary}\label{cor:1.21}
有双射
\[
  \Spa(E)^\diamond(R,R^+)\simeq\bigl\{\text{次数 }1\text{ 本原元}\in W_{\oo_E}(R^\circ)\bigr\}\big/W_{\oo_E}(R^\circ)^\times.
\]
\end{corollary}

现在，公式 $Y^\diamond_{S,E}=S\times\Spa(E)^\diamond$ 告诉我们：$S$ 的任何 untilt $S^\sharp$ 都定义一个态射 $S=S^{\sharp,\diamond}\to Y^\diamond_{S,E}$，也就是说一个态射 $S^\sharp\to Y_{S,E}$，而它事实上是局部由次数 $1$ 本原元定义的闭浸入。于是有：

\begin{proposition}\label{prop:1.22}
有双射
\[
\begin{aligned}
  &\{\text{$S$ 的 untilt}\}\\
  &\simeq\ \{\text{闭浸入}\ T\hookrightarrow Y_S,\ \text{局部由次数 }1\text{ 本原元定义}\}.
\end{aligned}
\]
\end{proposition}

此外可以验证：若 $S^\sharp$ 是 $S$ 的 untilt，则闭浸入 $S^\sharp\hookrightarrow Y_S$ 诱导闭浸入
\[
  S^\sharp\hookrightarrow X_S.
\]
Frobenius 在 $S$ 的 untilt 上有双射的作用，由公式 $(S^\sharp,\iota)\mapsto(S^\sharp,\iota\circ\mathrm{Frob})$ 给出。这由 $\varphi$ 在次数 $1$ 本原元集合上的作用给出。于是有下面的命题。

\begin{proposition}\label{prop:1.23}
有映射
\[
  \{\text{$S$ 的 untilt}\}/\mathrm{Frob}^\ZZ\ \longrightarrow\ \{\text{闭浸入}\ T\hookrightarrow X_S,\ \text{局部由次数 }1\text{ 本原元定义}\}.
\]
\end{proposition}

\subsubsection*{1.4.3. 若干 Hilbert 菱形}% 原文 1.4.3. Some Hilbert diamonds

最后说明：对任何 $d\ge1$，都有次数 $d$ 本原元的概念。确切地说，$f=\sum_{n\ge0}[a_n]\pi^n\in\mathbf{A}$ 称为\emph{次数 $d$ 本原的}，若 $a_0\in R^{\circ\circ}\cap R^\times$，$a_1,\dots,a_{d-1}\in R^{\circ\circ}$，且 $a_d\in R^{\circ\times}$。于是 $\mathbf{A}$ 中次数 $d$ 的本原元模 $\mathbf{A}^\times$ 与
\[
  \bigl[(\Spa(E)^\diamond)^d/\mathfrak{S}_d\bigr](R,R^+)
\]
一一对应，其中按对称群取商是对 [27] 中引入的忠实拓扑而言。因此，次数 $d$ 本原元与曲线的对称积相关：
\[
  X^{[d],\diamond}_S=(S\times\Spa(E)^{\diamond,d})/\varphi^\ZZ\times\mathfrak{S}_d.
\]
这些\emph{Hilbert 菱形}似乎相当有趣\yizhu{次数 $d$ 本原元对应 $(\Spa(E)^\diamond)^d/\mathfrak{S}_d$，即曲线的 $d$ 次对称积；「Hilbert 菱形」之名来自它与 Hilbert 概型的类比。}{A：原文 §1.4.3 p.10}，但我们不再多说，因为表述我们的猜想并不需要它们（尽管毫无疑问它们会出现在猜想的证明中）。
